\documentclass[11pt,a4paper]{article}

% ── Packages ──────────────────────────────────────────────────────────────────
\usepackage[top=2.8cm, bottom=2.8cm, left=3cm, right=3cm]{geometry}
\usepackage[T1]{fontenc}
\usepackage[utf8]{inputenc}
\usepackage{lmodern}
\usepackage{microtype}
\usepackage[colorlinks=true, linkcolor=linkblue, urlcolor=accent,
            pdftitle={L1-07 Caching -- System Design Foundations},
            bookmarks=true]{hyperref}
\usepackage{xcolor}
\usepackage{titlesec}
\usepackage{enumitem}
\usepackage{booktabs}
\usepackage{tabularx}
\usepackage{tcolorbox}
\usepackage{fancyhdr}
\usepackage{amsmath}
\usepackage{parskip}
\usepackage{mdframed}
\usepackage{graphicx}
\usepackage{tikz}
\usepackage{setspace}
\usepackage{soul}
\usepackage{array}
\usepackage{longtable}

% ── Colors ────────────────────────────────────────────────────────────────────
\definecolor{primary}{HTML}{1A1A2E}
\definecolor{accent}{HTML}{C0392B}
\definecolor{highlight}{HTML}{1A5276}
\definecolor{linkblue}{HTML}{154360}
\definecolor{lightbg}{HTML}{F4F6F7}
\definecolor{quizbg}{HTML}{EBF5FB}
\definecolor{termbg}{HTML}{EAF2FF}
\definecolor{curiobg}{HTML}{F5EEF8}
\definecolor{greenbg}{HTML}{EAFAF1}
\definecolor{notebg}{HTML}{FDF2E9}
\definecolor{accentlight}{HTML}{FADBD8}

% ── Section styling ───────────────────────────────────────────────────────────
\titleformat{\section}
  {\large\bfseries\color{highlight}}
  {\thesection.}{0.6em}{}
  [\vspace{2pt}\color{highlight!60}\rule{\linewidth}{0.6pt}]

\titleformat{\subsection}
  {\normalsize\bfseries\color{primary}}
  {\thesubsection}{0.5em}{}

\titleformat{\subsubsection}
  {\normalsize\itshape\color{primary!80}}
  {}{0em}{}

\titlespacing{\section}{0pt}{18pt}{8pt}
\titlespacing{\subsection}{0pt}{12pt}{4pt}
\titlespacing{\subsubsection}{0pt}{8pt}{2pt}

% ── Header/Footer ─────────────────────────────────────────────────────────────
\pagestyle{fancy}
\fancyhf{}
\fancyhead[L]{\small\color{highlight}\textbf{L1-07 --- Caching}}
\fancyhead[R]{\small\color{primary!70}System Design Interview Prep}
\fancyfoot[C]{\small\color{primary!60}\thepage}
\renewcommand{\headrulewidth}{0.4pt}
\renewcommand{\headrule}{\hbox to\headwidth{\color{highlight!40}\leaders\hrule height \headrulewidth\hfill}}

% ── Box environments ──────────────────────────────────────────────────────────
\tcbuselibrary{skins, breakable, hooks}

\newtcolorbox{termbox}[1]{
  enhanced, breakable,
  colback=termbg, colframe=highlight, coltitle=white,
  fonttitle=\bfseries\small, title={#1},
  left=8pt, right=8pt, top=6pt, bottom=6pt,
  attach boxed title to top left={yshift=-2.5mm, xshift=6mm},
  boxed title style={colback=highlight, sharp corners, left=4pt, right=4pt}
}

\newtcolorbox{industrybox}[1]{
  enhanced, breakable,
  colback=greenbg, colframe=highlight!50!black, coltitle=white,
  fonttitle=\bfseries\small, title={Industry Data: #1},
  left=8pt, right=8pt, top=6pt, bottom=6pt,
  attach boxed title to top left={yshift=-2.5mm, xshift=6mm},
  boxed title style={colback=highlight!60!black, sharp corners, left=4pt, right=4pt}
}

\newtcolorbox{trapbox}[1]{
  enhanced,
  colback=accentlight, colframe=accent, coltitle=white,
  fonttitle=\bfseries\small, title={Interview Trap: #1},
  left=8pt, right=8pt, top=6pt, bottom=6pt,
  attach boxed title to top left={yshift=-2.5mm, xshift=6mm},
  boxed title style={colback=accent, sharp corners, left=4pt, right=4pt}
}

\newtcolorbox{thinkbox}{
  enhanced,
  colback=notebg, colframe=accent!60,
  left=8pt, right=8pt, top=6pt, bottom=6pt,
  leftrule=3pt, rightrule=0pt, toprule=0pt, bottomrule=0pt,
}

\newtcolorbox{curiobox}[1]{
  enhanced, breakable,
  colback=curiobg, colframe=primary!60, coltitle=white,
  fonttitle=\bfseries\small, title={Q: #1},
  left=8pt, right=8pt, top=6pt, bottom=6pt,
  attach boxed title to top left={yshift=-2.5mm, xshift=6mm},
  boxed title style={colback=primary!70, sharp corners, left=4pt, right=4pt}
}

\newtcolorbox{quizbox}{
  enhanced, breakable,
  colback=quizbg, colframe=highlight, coltitle=white,
  fonttitle=\bfseries, title={Self-Quiz},
  left=10pt, right=10pt, top=8pt, bottom=8pt,
  attach boxed title to top left={yshift=-2.5mm, xshift=6mm},
  boxed title style={colback=highlight, sharp corners, left=4pt, right=4pt}
}

\newtcolorbox{answerbox}{
  enhanced, breakable,
  colback=lightbg, colframe=primary!40, coltitle=white,
  fonttitle=\bfseries\small, title={Model Answers},
  left=10pt, right=10pt, top=8pt, bottom=8pt,
  attach boxed title to top left={yshift=-2.5mm, xshift=6mm},
  boxed title style={colback=primary!60, sharp corners, left=4pt, right=4pt}
}

\newtcolorbox{insightbox}{
  enhanced,
  colback=primary!5, colframe=primary,
  left=8pt, right=8pt, top=6pt, bottom=6pt,
  leftrule=4pt, rightrule=0.4pt, toprule=0.4pt, bottomrule=0.4pt,
}

\setstretch{1.15}

% ── Document ──────────────────────────────────────────────────────────────────
\begin{document}

% ── Title page ────────────────────────────────────────────────────────────────
\begin{titlepage}
\vspace*{2cm}
\begin{center}
  {\color{primary}\rule{\linewidth}{3pt}}\\[12pt]
  {\huge\bfseries\color{highlight} L1-07}\\[6pt]
  {\Huge\bfseries\color{primary} Caching}\\[10pt]
  {\large\color{primary!70} Eviction Policies, Write Strategies, Redis Internals, and the Art of Avoiding Your Database}\\[20pt]
  {\color{primary}\rule{\linewidth}{1pt}}\\[16pt]
  {\normalsize\color{primary!70}
    \textit{Layer 1 --- Foundations} \quad\textbullet\quad
    \textit{System Design Interview Preparation} \quad\textbullet\quad
    \textit{2026 Edition}
  }\\[8pt]
  {\small\color{primary!60} Industry data sourced from Redis Labs, Facebook Engineering, Twitter/X Engineering, Cloudflare, DoorDash Engineering, and Slack Engineering blogs (2022--2025)}
\end{center}
\vfill
\begin{insightbox}
\textbf{Why this lesson exists.} Caching appears in every single system design interview without exception. The leaderboard question needs Redis sorted sets. The URL shortener needs a cache layer. The feed system needs result caching. The notification system you designed in session\_001 needs a device-token cache to avoid hammering the database on every fanout. If you design any system without discussing caching, the interviewer will ask: ``Where would you add a cache?'' You need to be ready not just with \emph{where} but \emph{which pattern}, \emph{which eviction policy}, and \emph{what happens when the cache fails}. This lesson covers all of it.
\end{insightbox}
\vspace{1cm}
\tableofcontents
\end{titlepage}

\newpage

% ── Section 1: Why This Exists ────────────────────────────────────────────────
\section{Why Caching Exists: The Latency and Load Problem}

In the early 2000s, every popular website ran on a roughly similar architecture: a web server that fielded requests and a relational database that stored everything. This worked fine at low traffic. When a page was requested, the server queried the database, assembled the HTML, and sent it back. The database could handle it. Everything was fine.

Then the traffic grew. Not gradually --- virally. Digg, the social bookmarking site, would occasionally surface a story that would send hundreds of thousands of visitors to a small website in the span of hours. This phenomenon was called the ``Digg Effect,'' and it killed servers. The same pattern replicated when Reddit sent traffic, when a story made the front page of Hacker News, when a tweet went viral. The common thread: a relational database, tuned for careful writes and complex queries, is terrible at handling ten thousand identical reads per second for the same data.

A relational database read involves parsing the query, planning an execution path, scanning indexes, reading data from disk (even with the buffer pool, cold data requires I/O), assembling result rows, and serializing them for the network. A simple \texttt{SELECT} that touches cached pages in the buffer pool takes 1--10 milliseconds. A query that misses the buffer pool and hits disk takes 10--100 milliseconds. At 10,000 requests per second, you need 10,000 database reads per second --- a load that most databases, absent extraordinary hardware, cannot sustain without becoming the bottleneck.

The insight that powered a generation of scalable web architecture was this: \emph{most data does not change frequently, but it is read constantly}. The 100 trending stories on the front page of Reddit are the same for every visitor for 60 seconds at a time. Computing that list requires a complex database query. But the result --- a ranked list of posts --- is stable between refreshes. If you compute it once and store the result somewhere fast, you can serve a thousand users from that one computation.

\begin{insightbox}
\textbf{The core principle of caching:} Trade memory for latency and database load. Memory is fast (100--300 nanoseconds for a cache read in the same data center) and relatively cheap. Database reads are slow (1--100 milliseconds) and expensive (each one consumes CPU, I/O, and a connection). Caching is the practice of keeping a small, fast copy of data close to where it is consumed, so that the slow, expensive original is consulted only when necessary.
\end{insightbox}

The numbers make the value concrete. Facebook's Memcached deployment, documented in their 2013 NSDI paper, stored hundreds of billions of objects in memory across their fleet. At peak, their cache reduced the number of database queries by more than 99\% for many workloads. That is not a typo: one in every hundred requests reached the database. The other 99 were served from memory. Without the cache, Facebook's database fleet would have needed to be roughly 100 times larger. The cache did not just make the system faster; it made it economically viable.

% ── Section 2: Cache-Aside and Write Patterns ──────────────────────────────────
\section{How Caches Interact with the Database: Write Patterns}

The trickiest part of caching is not reading --- it's writing. When data changes, the cache and the database can diverge. The four write patterns describe the four different answers to the question: when the application writes data, in what order does it write to the cache and to the database?

\subsection{Cache-Aside (Lazy Loading): The Most Common Pattern}

Cache-aside is the pattern you will use in the majority of system design interviews unless there is a specific reason to use something else. It is called ``lazy'' because the cache is populated only when data is first requested, not proactively.

The read path is: the application checks the cache for the requested key. If the data is there (a cache hit), return it immediately. If the data is not there (a cache miss), the application reads from the database, stores the result in the cache with a TTL (time-to-live), and returns the result. The write path is simpler: the application writes to the database, and either invalidates the cache key or updates the cache entry.

\begin{termbox}{Cache-Aside (Lazy Loading)}
The application owns the cache interaction. On a read miss, the application fetches from the database and populates the cache. On a write, the application writes to the database and invalidates or updates the corresponding cache entry. The cache is never written without a preceding read miss triggering it.
\end{termbox}

Cache-aside has three important properties. First, only data that is actually requested ends up in the cache. You do not waste memory caching data that nobody reads. Second, the cache can survive a failure gracefully: if the cache goes down, the application simply reads from the database on every request. It's slower, but it works. Third, after a restart or a new cache node joins the cluster, the cache starts cold (empty). The first wave of requests all miss, creating a spike of database load. This is called the cold-start problem, and it is mitigated by cache warming (covered later in this lesson).

\begin{thinkbox}
\textbf{Think about this:} In cache-aside, the application writes to the database first, then invalidates the cache. Why not write to the cache first? If you write to the cache and then the database write fails, you now have a cache with data that does not match the database. The cache will serve incorrect data until the TTL expires. By writing to the database first and then invalidating, a failed invalidation means the cache still has old data, but the database has correct data --- a much safer inconsistency. Old data expires eventually; corrupted data can persist indefinitely.
\end{thinkbox}

\subsection{Write-Through: Consistency at the Cost of Write Latency}

In a write-through cache, every write goes through the cache \emph{and} to the database, synchronously. When your application writes a value, the cache library (or a middleware layer) writes it to the cache and waits for the database write to complete before acknowledging success. Both stores are always updated together.

The advantage is consistency: the cache is never stale, because every write keeps it current. The disadvantage is write latency: every write pays the cost of \emph{both} a cache write and a database write before returning. If the database is the slow path (10ms), write-through doubles the latency of every write operation.

Write-through also causes the cache to fill with data that may never be read. If a user updates a profile field that is only read once a month, that write-through update writes the profile to the cache at the cost of a cache entry that will expire unused. This is the opposite of cache-aside's efficiency: cache-aside caches only what is read; write-through caches everything that is written.

\begin{industrybox}{Write-Through in Banking Systems}
Financial systems and payment processors frequently use write-through caching because they prioritize data consistency over write performance. A credit card authorization that updates account state must be immediately consistent: the next authorization check must see the up-to-date balance. In 2024, Stripe's engineering blog described their approach to their idempotency layer, which uses write-through semantics: every idempotency key write is committed synchronously to both their Redis cache and PostgreSQL before returning to the caller. The overhead is acceptable because payment operations are low-volume and high-value, not high-frequency.
\end{industrybox}

\subsection{Write-Back (Write-Behind): Fastest Writes, Real Risk}

Write-back caching inverts the priority: the application writes to the cache only, acknowledges success immediately, and the cache system asynchronously flushes the data to the database later --- in bulk, at intervals, or on eviction.

This is the fastest write pattern by a significant margin. The application's write latency is now a cache write (sub-millisecond), not a database write (1--10ms). The database is updated asynchronously, and because multiple writes to the same key can be coalesced, the database sees fewer operations: ten updates to the same user record in ten seconds may result in a single database write.

The risk is real: any data written to the cache but not yet flushed to the database is lost if the cache node fails. For a Redis write-back layer, this means data written in the last flush interval (often 1--5 seconds) is potentially lost. This is an acceptable tradeoff for session data, analytics counters, and other workloads where occasional data loss is tolerable. It is completely unacceptable for financial transactions, order records, or anything where data loss constitutes a legal or business failure.

\subsection{Write-Around: Skip the Cache on Write}

Write-around is a pattern where the application writes directly to the database, bypassing the cache entirely. The cache is populated only when that data is subsequently read. This sounds like cache-aside, but the difference is intentional: write-around is used specifically for data that is written frequently but read infrequently.

Log ingestion pipelines are a canonical example. You write hundreds of thousands of log entries per second to permanent storage. Each log entry will almost certainly never be read --- only a tiny fraction of logs are ever queried in investigations. Caching every log write would fill your cache with data that never gets a cache hit. Write-around lets you batch-write to the database efficiently without polluting the cache.

\begin{tabularx}{\linewidth}{lXXX}
\toprule
\textbf{Pattern} & \textbf{Read path} & \textbf{Write path} & \textbf{Best for} \\
\midrule
Cache-aside & Check cache, miss $\to$ DB, populate cache & Write DB, invalidate cache & General purpose; most common \\
Write-through & Always from cache (always current) & Write cache + DB synchronously & High read rate, consistency required \\
Write-back & Always from cache & Write cache only; async DB flush & High write rate, some loss tolerable \\
Write-around & Check cache, miss $\to$ DB & Write DB only, skip cache & Write-heavy, infrequently read data \\
\bottomrule
\end{tabularx}

% ── Section 3: Eviction Policies ───────────────────────────────────────────────
\section{Eviction Policies: What Leaves When Memory Fills Up}

A cache has finite memory. When it fills up, something must be evicted to make room for new data. The eviction policy determines \emph{which} entries are removed. Choosing the wrong policy for your access pattern can dramatically reduce cache effectiveness.

\subsection{LRU: Least Recently Used}

LRU evicts the entry that was accessed least recently. The idea is temporal locality: if you accessed something recently, you are likely to access it again soon. If you have not accessed something in a long time, you probably do not need it. LRU is implemented using a combination of a hash map (for O(1) lookup) and a doubly linked list (for O(1) promotion of recently-used entries to the front and eviction from the back).

LRU is the most commonly correct choice for web application caches. User request patterns exhibit strong recency bias: the articles read in the last hour are much more likely to be read again than articles read six months ago. Redis's default eviction policy is \texttt{allkeys-lru}, which evicts the least recently used key across all keys when memory is full.

\subsection{LFU: Least Frequently Used}

LFU evicts the entry that has been accessed the fewest times in total. The idea is frequency of access: if an entry has been requested a million times, it is probably important; if it has been requested once, it is probably transient. LFU is a better policy when access frequency is the relevant signal --- for example, in a media streaming platform where certain songs or videos are perennially popular across very long time windows (months, years).

The implementation challenge of LFU is that it must track access counts, which adds memory overhead and requires periodic decay (otherwise an item that was popular a year ago but is no longer requested retains a high count forever). Redis's \texttt{allkeys-lfu} policy uses an approximate LFU with logarithmic counter decay to handle this.

\subsection{TTL: Time-Based Expiry}

TTL (time-to-live) is not an eviction policy in the strict sense --- it is an expiration mechanism. Each cache entry is assigned a maximum lifetime. When the TTL expires, the entry is considered stale and will be evicted (lazily, when next accessed, or eagerly by a background process). TTL is used in conjunction with LRU or LFU, not as an alternative.

Setting the right TTL requires understanding the data's update frequency and the acceptable staleness window. A user's profile photo changes infrequently; a TTL of one hour is fine. A stock price changes every second; a TTL of one second may still be too long. A trending topics list changes every few minutes; a TTL of 30 seconds is a reasonable balance between freshness and database load.

\begin{thinkbox}
\textbf{Think about this:} If LRU is the most common choice, why does LFU exist? LRU has a known failure mode called ``cache pollution by scan.'' If your application performs a full table scan or iterates over millions of unique records in sequence, every record is accessed exactly once and then never again. Under LRU, these scan accesses promote their entries to the front of the recency list, evicting genuinely hot data that was serving many users. LFU handles this better: items accessed once do not accumulate frequency and are quickly evicted when memory pressure appears. Database buffer pools (like MySQL's InnoDB buffer pool) use variants of this insight to protect against scan pollution.
\end{thinkbox}

\subsection{FIFO and Random Eviction}

FIFO (First In, First Out) evicts the oldest entries by insertion order, regardless of access pattern. It is simple to implement but generally performs poorly: the first item inserted is not necessarily the least valuable. Random eviction picks a random entry to evict. Surprisingly, random eviction can perform comparably to LRU in some workloads (particularly Zipfian distributions where popular items are very popular and will be re-fetched quickly anyway). Redis uses random sampling internally as part of its approximate LRU implementation for performance reasons.

% ── Section 4: Redis Architecture ─────────────────────────────────────────────
\section{Redis: Architecture, Data Structures, and Why It Matters}

Redis (Remote Dictionary Server) is the dominant in-memory data store used for caching in production systems. Understanding why requires understanding both what Redis is and what it is not.

\subsection{The Single-Threaded Architecture and Why It Works}

Redis processes commands in a single thread. This is a deliberate design decision that dates to Redis's original architecture in 2009 by Salvatore Sanfilippo. The choice seems counterintuitive in a world of multi-core servers, but it has a profound advantage: because Redis never runs two commands concurrently, it never needs locks, mutexes, or any form of concurrency control. Every command is atomic relative to every other command.

The single thread is fast enough because Redis's bottleneck is almost never CPU: it is network I/O and memory bandwidth. The Redis event loop uses non-blocking I/O (via \texttt{epoll} on Linux) to multiplex thousands of client connections on a single thread. A modern Redis instance can handle 100,000--1,000,000 operations per second on a single core, with p99 latency under 1 millisecond.

Redis 6.0 (2020) introduced threaded I/O for reading and writing network sockets, while keeping command processing single-threaded. Redis 7.0 (2022) pushed further with background threads for persistence. The community project Dragonfly (2022) takes a fully multi-threaded approach, claiming 25x better throughput on multi-core hardware in benchmarks, though real-world Redis workloads at typical scale rarely exhaust single-thread capacity.

\begin{industrybox}{Redis at Twitter/X Scale (2023)}
Twitter's engineering team documented their Redis usage in 2023: they operate over 3,000 Redis instances holding approximately 70 terabytes of cached data. The primary use case is timeline caching: a user's home timeline is materialized and cached as a Redis list containing tweet IDs. When a user with 100,000 followers tweets, the system (previously called Flock, the fanout service) writes the tweet ID to up to 100,000 Redis lists in parallel. Timeline reads serve from this precomputed Redis list with sub-millisecond latency rather than recomputing from the database. This write-fanout model trades write amplification for read speed --- a classic caching tradeoff.
\end{industrybox}

\subsection{Redis Data Structures: Beyond Simple Key-Value}

The most common interview mistake about Redis is treating it as a simple key-value store equivalent to a dictionary. Redis is a data structure server. It stores complex types, each with their own operations that run server-side, atomically.

\textbf{Strings} are the simplest type: a key maps to a single value (which can be a string, integer, or binary blob up to 512 MB). \texttt{SET key value EX 3600} sets a key with a 1-hour TTL. \texttt{INCR counter} atomically increments an integer (no read-modify-write race conditions). Strings are used for caching serialized objects, distributed locks, and rate-limiting counters.

\textbf{Hashes} store a map of field-value pairs under a single key. \texttt{HSET user:123 name ``Alice'' email ``alice@example.com''} stores a user record as a hash. \texttt{HGET user:123 name} retrieves just the name field without deserializing the full object. Hashes are ideal for storing structured objects where you need field-level access.

\textbf{Lists} are linked lists of strings. \texttt{LPUSH} and \texttt{RPUSH} add to the left or right. \texttt{LRANGE key 0 49} retrieves the first 50 items. Lists model queues (\texttt{RPUSH} to enqueue, \texttt{BLPOP} to dequeue with blocking) and timelines (Twitter's list-per-user approach). The maximum size is $2^{32} - 1$ elements.

\textbf{Sets} are unordered collections of unique strings. \texttt{SADD}, \texttt{SREM}, and \texttt{SMEMBERS} manage membership. \texttt{SUNION}, \texttt{SINTER}, and \texttt{SDIFF} operate on multiple sets. A practical use: storing the set of users who liked a post (\texttt{SADD post:456:likes user:789}), then checking (\texttt{SISMEMBER}) with O(1) complexity.

\textbf{Sorted sets (ZSets)} are the data structure that most frequently appears in system design interviews. A sorted set is a set of members, each associated with a floating-point score. The set is maintained in sorted order by score at all times. \texttt{ZADD leaderboard 9823.5 user:42} adds user 42 with score 9823.5. \texttt{ZRANGE leaderboard 0 9 REV WITHSCORES} returns the top 10 entries in descending order. \texttt{ZRANK leaderboard user:42} returns the rank (position) of user 42 in O(log N).

\begin{termbox}{Redis Sorted Set --- The Leaderboard Structure}
A ZSet stores (member, score) pairs, maintaining members in sorted score order. Core operations: \texttt{ZADD} (insert/update, O(log N)), \texttt{ZRANGE} with \texttt{REV} flag (top-N retrieval, O(log N + M)), \texttt{ZRANK} (rank of a member, O(log N)), \texttt{ZINCRBY} (increment a member's score atomically, O(log N)). Because all operations are O(log N) server-side, a leaderboard with 10 million entries is queried as fast as one with 1,000 entries from the application's perspective --- network round-trip time dominates.
\end{termbox}

\textbf{Bitmaps} treat a string as a bit array. \texttt{SETBIT dau:20250101 user\_id 1} marks a user as active on a specific day. \texttt{BITCOUNT dau:20250101} counts the number of active users. A bitmap for 100 million users requires only 12.5 MB. This is how systems track daily active users (DAU) at massive scale without per-user database writes.

\textbf{HyperLogLog} is a probabilistic data structure that estimates the cardinality of a set using only 12 KB of memory, with at most 0.81\% error. \texttt{PFADD visitors:page-X user-id} adds a visitor. \texttt{PFCOUNT visitors:page-X} returns the estimated unique visitor count. The 12 KB limit is constant regardless of how many unique elements have been added. Used for unique visitor counting, unique query counting, and other cardinality estimation problems where approximate counts are acceptable.

\begin{thinkbox}
\textbf{Think about this:} Why would you use HyperLogLog instead of a Set for unique visitor counting? A Redis Set of 10 million unique user IDs requires roughly 500 MB of memory (50 bytes per ID average). A HyperLogLog for 10 million IDs requires exactly 12 KB. The trade: you lose the ability to ask ``did this specific user visit?'' (the Set supports that; HyperLogLog does not) but you save 99.997\% of the memory. For dashboards where you just need the count, not the membership query, HyperLogLog is the right tool.
\end{thinkbox}

\subsection{Redis Persistence: RDB and AOF}

Redis is an in-memory store, which means a server restart loses all data unless persistence is configured. Redis provides two mechanisms.

\textbf{RDB (Redis Database)} snapshots are point-in-time dumps of the entire dataset to disk, written periodically (e.g., every 15 minutes if at least one key changed). RDB files are compact and fast to restore. The risk: data written since the last snapshot is lost on failure. For a cache layer where losing 15 minutes of cached data simply means more cache misses (not data corruption), RDB is usually acceptable.

\textbf{AOF (Append-Only File)} logs every write command. On restart, Redis replays the log to reconstruct the dataset. AOF can be configured to fsync every command (no data loss, but slower writes), every second (at most 1 second of data loss), or never (OS decides when to flush). AOF provides stronger durability guarantees at the cost of larger files and slower restores.

For a pure caching layer, many teams disable persistence entirely: the cache is considered ephemeral, and a cold restart is handled by the cache warming process rather than data recovery.

\subsection{Redis Cluster: Horizontal Scaling}

A single Redis node is limited to the memory of one machine (practical limit: 32--64 GB for reasonable performance). Redis Cluster shards data across multiple nodes using consistent hashing on key slots. There are 16,384 hash slots total. Each key is assigned to a slot (\texttt{CRC16(key) mod 16384}), and each slot is owned by one primary node (with one or more replicas).

Redis Cluster provides automatic failover: if a primary node fails, its replica is promoted within seconds (typically under 10 seconds). The cluster requires a minimum of 3 primary nodes to achieve quorum for leader election. In practice, production clusters run 6 or more nodes (3 primaries + 3 replicas) for meaningful redundancy.

The limitation of Redis Cluster is that multi-key commands (like \texttt{MGET}) and transactions only work correctly if all involved keys are on the same node. Keys can be forced into the same slot using hash tags: \texttt{\{user:123\}:profile} and \texttt{\{user:123\}:settings} will always be on the same node because Redis hashes only the portion within the curly braces.

\subsection{Redis Sentinel: High Availability Without Sharding}

Redis Sentinel is the alternative to Redis Cluster for high-availability when you do not need horizontal scaling. Sentinel runs as a separate process that monitors Redis master and replica nodes. If the master becomes unreachable, Sentinel nodes reach a quorum vote and automatically promote a replica to master. Client libraries that speak the Sentinel protocol automatically discover the new master address. Sentinel requires a minimum of 3 Sentinel processes to prevent split-brain election.

\begin{tabularx}{\linewidth}{lXX}
\toprule
 & \textbf{Redis Sentinel} & \textbf{Redis Cluster} \\
\midrule
Primary use case & High availability & Horizontal scaling \\
Data partitioning & No (all data on one master) & Yes (16,384 hash slots) \\
Multi-key ops & Fully supported & Limited (same slot only) \\
Minimum nodes & 1 master + 2 replicas + 3 Sentinels & 3 primaries + 3 replicas \\
Failover time & 10--30 seconds & 10--15 seconds \\
Use when & Data fits one machine, need HA & Data exceeds one machine \\
\bottomrule
\end{tabularx}

% ── Section 5: Redis vs Memcached ─────────────────────────────────────────────
\section{Redis vs Memcached: The Right Tool for the Job}

Memcached is the cache system that predates Redis in the industry, and the Redis vs. Memcached debate appears in interviews regularly. The answer today is nuanced but trends strongly toward Redis.

Memcached is a pure key-value cache. Every value is a blob (bytes). There are no data structures, no persistence, no pub/sub, no scripting. Memcached is designed for one thing and does it well: store arbitrary blobs by key, retrieve them fast, evict them when memory fills up. Its multi-threaded architecture makes it more efficient than Redis on multi-core servers for pure key-value workloads.

Facebook built the most famous Memcached deployment in the world: their ``Scaling Memcached at Facebook'' paper (NSDI 2013) describes a fleet serving 1 billion requests per second at peak, with a custom routing layer called McRouter that aggregated replication and failover logic. This paper is worth reading once; interviewers occasionally reference it.

\begin{industrybox}{Facebook Memcached at Scale (NSDI 2013)}
At the time of the paper, Facebook operated over 800 servers in their largest Memcached cluster, storing over 1,000 terabytes of data across their fleet. McRouter, their Memcached proxy, handled 1 billion requests per second at peak. The key insight from the paper: at Facebook's scale, the bottleneck was not Memcached's performance but the \emph{coordination} overhead --- ensuring consistent invalidation across replicated clusters when data changed. They invented a lease mechanism to handle cache stampede (detailed in Section 6) and a concept called ``regional pools'' to handle data that needs different replication semantics.
\end{industrybox}

The modern recommendation (AWS, Google Cloud, Redis Labs, and most engineering blogs post-2020) is to default to Redis. The reasons are: Redis's data structures reduce the need for complex application-side logic (a sorted leaderboard in Redis is two commands; in Memcached it requires reading the entire leaderboard, sorting it in the application, and writing it back atomically using CAS); Redis's persistence options provide durability when needed; Redis's pub/sub and streams capabilities extend it beyond pure caching; and Redis's active development community means better operational tooling. Memcached remains a valid choice when you have an existing Memcached deployment or when the specific multi-threaded performance of Memcached on CPU-bound key-value workloads is the critical requirement.

% ── Section 6: Cache Stampede and Thundering Herd ─────────────────────────────
\section{Cache Stampede and the Thundering Herd Problem}

Cache stampede is one of the most dangerous failure modes in caching systems, and it is a guaranteed interview topic at senior level. Understanding it precisely and being able to name the solutions is essential.

\subsection{The Anatomy of a Stampede}

Imagine your application caches the result of a popular, expensive database query with a TTL of 60 seconds. The key expires. In the millisecond window between expiry and repopulation, 500 concurrent requests arrive for that key. All 500 experience a cache miss simultaneously. All 500 query the database. The database, sized to handle normal load with caching, is now receiving 500 times its expected query rate for this resource. It slows down. This causes other queries to back up. Other cache TTLs expire while the database is overloaded. The cascade accelerates. This is a stampede.

The thundering herd is the same phenomenon at a larger scale: a cache node restart brings up with a cold (empty) cache, and every key that existed in that node now misses simultaneously, sending a thundering herd of queries to the backend all at once.

\begin{trapbox}{Assuming Cache Failures Are Rare}
In interviews, candidates often design cache layers without discussing what happens when cache nodes fail, restart, or TTLs expire simultaneously. This is a red flag. Interviewers at senior level will always probe: ``What happens if your cache is empty?'' You must be ready to explain cache stampede and describe at least one mitigation. The answer ``the database will handle it'' is always wrong --- if the database could handle the full load without caching, you would not need the cache.
\end{trapbox}

\subsection{Solutions to Cache Stampede}

\textbf{Mutex / lock on cache miss.} When a cache miss occurs, before querying the database, the application tries to acquire a distributed lock for that cache key. Only the first requester acquires the lock and proceeds to the database. All other requesters wait for the lock to be released, at which point the first requester has repopulated the cache. They re-check the cache and now find the data. This prevents parallel database queries for the same key. The implementation in Redis uses \texttt{SET key NX PX 5000} (set if not exists, with a 5-second expiry as a safety timeout).

\textbf{Probabilistic early expiration (PER).} Instead of waiting for the TTL to reach exactly zero, each reader independently uses a formula that probabilistically triggers a cache refresh \emph{before} the TTL expires. The closer to expiry the entry is, the higher the probability of triggering a refresh on any individual read. This is implemented as: if \texttt{current\_time - delta * beta * log(random())} exceeds the item's expiry time, recompute. The logarithm ensures that as the expiry approaches, increasingly many reads trigger the refresh, so the cache is almost certainly refreshed before it expires for real. No lock is needed.

\textbf{Stale-while-revalidate.} The cache continues serving the stale (expired) value while a background process asynchronously recomputes the fresh value. The user sees slightly old data for one request cycle, but the database receives only one query instead of thousands. This is the correct pattern when serving stale data for a few hundred milliseconds is acceptable. HTTP CDNs implement this natively via the \texttt{Cache-Control: stale-while-revalidate=60} directive.

\textbf{Request coalescing.} Multiple identical in-flight requests for the same key are deduplicated at a proxy or middleware layer. Only one upstream request is issued; all coalesced requesters receive the same response when it returns. Varnish, Nginx's proxy cache, and Cloudflare's cache implement this via their ``request collapsing'' feature.

\begin{industrybox}{Facebook's Lease Mechanism (NSDI 2013)}
Facebook's NSDI 2013 paper introduced the ``lease'' mechanism specifically to solve cache stampede. When a key is first requested after a miss, Memcached returns a lease token (a unique 64-bit integer) along with the miss response. Only the holder of the lease is allowed to compute the new value and write it back. Other clients that request the same key while the lease is held receive a ``wait'' signal and retry after a short delay (typically 10ms). The lease expires automatically if the holder crashes. This is mathematically equivalent to the mutex approach but integrated into the cache protocol itself, eliminating the need for a separate distributed lock service.
\end{industrybox}

% ── Section 7: Cache Warming ──────────────────────────────────────────────────
\section{Cache Warming: Avoiding the Cold Start}

A cold cache --- whether from a new deployment, a node restart, or scaling out new cache nodes --- is a performance cliff. The first wave of requests all miss, the database receives the full load it was shielded from, latency spikes, and in extreme cases the database becomes the bottleneck and causes a cascade. Cache warming is the practice of pre-populating the cache before it serves traffic.

\subsection{Warming Strategies}

\textbf{Pre-population at startup.} Before bringing a new cache node online, run a job that reads the most-accessed keys from the database (identified via query logs, analytics, or a recent key dump) and populates the cache. This requires knowing which keys are hot, which can be learned from metrics on the existing cache nodes. Cloudflare publishes their approach: when a new edge cache node is provisioned, it immediately subscribes to a peer-to-peer cache warming protocol that pulls hot objects from neighboring nodes rather than from the origin server.

\textbf{Shadow traffic replay.} Record real production requests for a time window and replay them against the new cache. This is the most accurate warming strategy because it reproduces the actual access pattern. The replay can run at an accelerated rate (10x or 100x) to warm the cache quickly.

\textbf{Gradual traffic migration.} Instead of cutting all traffic to a new cache at once, route 1\%, then 5\%, then 20\%, then 100\% over time. The cache warms organically under real traffic while the database absorbs the elevated miss rate during the ramp-up. This is the simplest strategy and the most commonly used in practice because it requires no special tooling.

\textbf{Cache-aside with negative TTLs.} For critical hot keys (homepage data, trending content, price lists), seed them explicitly in the startup sequence and set a longer TTL so they do not expire at the same time. Stagger TTLs by adding a random jitter of ±10\% to prevent synchronized mass expiry across all keys.

\begin{industrybox}{DoorDash Cache Warming (2023)}
DoorDash's engineering blog documented their cache warming approach for their restaurant menu cache in 2023. Menus are cached with 15-minute TTLs. When deploying a new cache cluster, they run a warming job that reads the 50,000 most-frequently-accessed restaurant menu keys (identified from the previous cluster's hit logs) and pre-populates them. The warming job completes in under 3 minutes. Without it, the first 15 minutes after a new cluster comes online sees a 40--60\% cache miss rate (compared to steady-state 2--5\%). With warming, the initial miss rate is under 10\%.
\end{industrybox}

% ── Section 8: Distributed Caching ────────────────────────────────────────────
\section{Distributed Caching: Consistent Hashing and Sharding}

When your data outgrows a single cache node, you must distribute the cache across multiple nodes. The central problem is: given a cache key, which node holds it? And what happens when a node is added or removed?

\subsection{Naive Hashing and Its Failure Mode}

The naive approach is modulo hashing: node = hash(key) mod N, where N is the number of nodes. This works when N is fixed. When a node fails and N becomes N-1, every key's node assignment changes. If you had 10 nodes and one fails, you go from hash(key) mod 10 to hash(key) mod 9. Almost every key now maps to a different node. Every cached value is effectively lost. Your entire cache goes cold simultaneously.

\subsection{Consistent Hashing}

Consistent hashing solves this by mapping both keys and nodes to positions on a conceptual ring (a circle of integers 0 to $2^{32}$). Each key is assigned to the first node whose position on the ring is equal to or greater than the key's hash (clockwise traversal). When a node is added or removed, only the keys that were assigned to that node need to be remapped. If you have 10 nodes and remove one, only $\frac{1}{10}$ of all keys are remapped. The other $\frac{9}{10}$ continue to resolve to the same nodes.

\begin{termbox}{Consistent Hashing}
A key distribution technique where both keys and nodes are hashed to positions on a conceptual ring. A key is owned by the nearest node clockwise on the ring. Adding or removing a node remaps only keys previously owned by adjacent nodes, minimizing cache disruption. Virtual nodes (vnodes) address the imbalance problem: each physical node is represented by multiple positions on the ring, so removal of one node distributes its keys across many other nodes rather than overloading one neighbor.
\end{termbox}

Virtual nodes (vnodes) solve the load imbalance problem in consistent hashing. If you place each physical node at one position on the ring, when you remove a node, all its keys go to a single neighbor. With 100--200 vnodes per physical node, removing one physical node distributes its keys across dozens of physical nodes proportionally.

\subsection{Consistent Hashing in Practice}

Redis Cluster uses a slightly different approach: 16,384 fixed hash slots rather than a continuous ring. The division into 16,384 slots was chosen to fit the cluster's configuration (which node owns which slots) into a 2 KB bitmap. The behavior is functionally equivalent to consistent hashing for most purposes. Each primary node owns a contiguous range of slots, and moving a shard means reassigning a range of slots from one node to another.

Memcached clients (like the libmemcached library) implement consistent hashing natively to distribute keys across a Memcached cluster. When a node is added, only the keys in the ring segment adjacent to the new node are affected.

% ── Section 9: Decision Framework ─────────────────────────────────────────────
\section{Decision Framework: Choosing the Right Cache Strategy}

\subsection{Choosing a Write Pattern}

Start by asking: what does data consistency cost in your domain?

If you are caching database query results, user sessions, or computed aggregates where serving slightly stale data for a TTL window is acceptable, use \textbf{cache-aside}. It is the default. It is the most operationally simple. It handles cache failures gracefully. It does not cache data that is never read.

If you have a high read-to-write ratio and the cost of a cache miss (going to the database) is too high to accept even during TTL expiry, use \textbf{write-through}. The cache is always warm. The trade is write latency. This is appropriate for session stores and user preference caches where every user interaction involves a read.

If your workload is write-heavy with bursty patterns (IoT sensor data, analytics event counters, real-time score updates) and you can tolerate a small window of data loss, use \textbf{write-back} with Redis AOF at the 1-second fsync interval as a safety net.

If your cache serves hot reads but writes are infrequent or the written data is rarely read back, use \textbf{write-around} to avoid polluting the cache.

\subsection{Choosing an Eviction Policy}

If your access pattern is dominated by recency (news feeds, notification histories, recently viewed products), use \textbf{LRU}. This is the right default for most web applications.

If your access pattern has long-lived popularity (music libraries, product catalogs, reference data), use \textbf{LFU}. Items that are popular over weeks will not be evicted by a spike of one-time accesses.

Always combine your eviction policy with \textbf{TTL} to ensure stale data does not accumulate indefinitely. Never rely on eviction alone to remove outdated data.

\subsection{Choosing Between Redis Sentinel, Cluster, or Single Node}

If your data fits in 16--32 GB and you need high availability (automatic failover), use \textbf{Redis Sentinel}. If your data exceeds one machine or your throughput requires multiple nodes, use \textbf{Redis Cluster}. If you are prototyping or the data is small and loss is acceptable (pure caching with no state), a \textbf{single Redis node} is fine.

% ── Section 10: Critical Numbers ─────────────────────────────────────────────
\section{Critical Numbers: What to Memorize}

\begin{tabularx}{\linewidth}{Xr}
\toprule
\textbf{Fact} & \textbf{Value} \\
\midrule
Redis single-node operations per second & 100K -- 1M ops/sec \\
Redis p99 read latency (same data center) & $<$ 1 ms \\
Redis p99 read latency (cross-region) & 1 -- 5 ms \\
L1 CPU cache access & $\sim$1 ns \\
RAM access & $\sim$100 ns \\
SSD random read & $\sim$100 $\mu$s \\
HDD random read & $\sim$10 ms \\
Database query (cached pages, no disk I/O) & 1 -- 10 ms \\
Database query (cold, disk I/O) & 10 -- 100 ms \\
Redis max memory per instance (practical) & 32 -- 64 GB \\
Redis hash slots (Cluster total) & 16,384 \\
Redis HyperLogLog memory usage & 12 KB (constant) \\
Redis bitmap for 100M users & 12.5 MB \\
Facebook Memcached peak throughput (2013) & 1 billion reqs/sec \\
Twitter Redis instances (2023) & $>$3,000 nodes, $\sim$70 TB \\
Good cache hit ratio target (web apps) & $>$ 95\% \\
Acceptable hit ratio (aggressive TTLs) & 80 -- 90\% \\
Consistent hashing: keys remapped on node removal & $\frac{1}{N}$ of all keys \\
Naive mod hashing: keys remapped on node removal & $\frac{N-1}{N}$ of all keys \\
\bottomrule
\end{tabularx}

% ── Section 11: System Design Applications ────────────────────────────────────
\section{System Design Application: Three Interview Examples}

\subsection{Leaderboard with Redis Sorted Sets}

The question: design a real-time gaming leaderboard for 10 million players, updated continuously, that can serve rank queries (``what is my rank?'') and top-N queries (``show me the top 100 players'').

The naive database approach is a \texttt{SELECT rank FROM players ORDER BY score DESC LIMIT 100} query. At 10 million rows, this query requires a full index scan and sorting. At high concurrency, this becomes a bottleneck quickly.

The Redis approach: maintain a sorted set with key \texttt{leaderboard:global}. When a player's score changes, call \texttt{ZADD leaderboard:global newscore playerid} (O(log N)). To get the top 100: \texttt{ZRANGE leaderboard:global 0 99 REV WITHSCORES} (O(log N + 100)). To get a specific player's rank: \texttt{ZREVRANK leaderboard:global playerid} (O(log N)). All operations are sub-millisecond, regardless of the 10 million player count. The sorted set automatically maintains sorted order on every update.

For regional leaderboards, use separate keys (\texttt{leaderboard:NA}, \texttt{leaderboard:EU}). For time-windowed leaderboards (weekly, monthly), use a key per window and delete old windows on expiry. For daily leaderboards, use \texttt{leaderboard:2025-01-15} and set a TTL of 48 hours so historical leaderboards auto-expire.

\begin{industrybox}{Riot Games Leaderboard (2022)}
Riot Games (League of Legends, Valorant) documented their ranked leaderboard system in 2022. They use Redis sorted sets to maintain per-region, per-queue leaderboards. The Challenger tier leaderboard (top 300 players per region) is served entirely from Redis with no database hit on reads. Score updates (rank changes after a game) are written to the sorted set immediately and persisted to the database asynchronously. The leaderboard page handles peak loads of over 500,000 reads per minute with sub-5ms p99 latency.
\end{industrybox}

\subsection{Session Store}

The question: how would you store authenticated user sessions for a web application serving 50 million daily active users?

Sessions are accessed on every API request. With 50 million DAU and an average of 20 requests per session, that is 1 billion session lookups per day, or roughly 12,000 reads per second. A session lookup from a relational database takes 1--10ms. At 12,000 reads per second, the database would spend most of its time serving session queries, leaving little capacity for business logic queries.

The solution: Redis as a session store. Use \texttt{SET session:\{session\_id\} \{serialized\_session\} EX 86400} (24-hour TTL). Reads are \texttt{GET session:\{session\_id\}} at sub-millisecond latency. Writes are \texttt{SET} on login, \texttt{DEL} on logout, \texttt{EXPIRE} to extend on activity.

For horizontal scaling of the session store, use write-through semantics: every session write goes to both Redis and the database (for durability and offline processing). Reads come exclusively from Redis. On a cache miss (Redis node failure or eviction), fall back to the database. This is the pattern used by AWS ElastiCache in session management documentation.

\subsection{Database Query Result Cache}

The question: a product detail page fetches product information, pricing, inventory, and reviews from the database. The page is loaded 100,000 times per day for popular products. How do you cache this?

Use cache-aside at the application layer. On the first request for product 456, the application queries the database (4 queries: product, pricing, inventory, reviews), serializes the result to JSON, and stores it as \texttt{SET product:456:full\_page \{json\} EX 300} (5-minute TTL). Subsequent requests for 300 seconds get the cached JSON and bypass all 4 database queries.

When a product is updated (price change, inventory update), the application writes to the database and then calls \texttt{DEL product:456:full\_page} to invalidate the cache. The next request repopulates it. This ensures stale data does not persist beyond one update cycle.

For a product catalog of 10 million items, storing full page caches requires significant memory: if an average product page serializes to 2 KB, 10 million cached products require 20 GB. Not every product needs caching --- apply cache-aside only to products that actually receive traffic. If 1\% of products drive 99\% of traffic (Zipf distribution, which is typical), caching 100,000 products covers 99\% of requests with 200 MB of cache memory.

% ── Curiosity Corner ──────────────────────────────────────────────────────────
\newpage
\section*{Curiosity Corner}
\addcontentsline{toc}{section}{Curiosity Corner}

\textit{These are the questions that likely occurred to you while reading. Take a moment to try answering each one before reading the explanation.}

\vspace{8pt}

\begin{curiobox}{What is Dragonfly DB, and is it replacing Redis?}
Dragonfly is an in-memory data store released in 2022 (dragonfly.io), designed as a Redis-compatible drop-in replacement with a fully multi-threaded architecture. It claims 25x higher throughput than Redis on the same hardware by using the ``Shared-nothing'' approach with per-thread data shards and a lock-free algorithm for managing data structures. In benchmarks published by the Dragonfly team, a single Dragonfly instance matches the throughput of a multi-node Redis Cluster. As of 2024, Dragonfly is in production at some companies but has not achieved the ecosystem maturity of Redis. The Redis ecosystem includes thousands of operator tools, monitoring integrations, client libraries, and patterns. Dragonfly is worth watching but is not yet the default recommendation. If your workload is bottlenecked on Redis single-thread CPU capacity (extremely rare at most scales), Dragonfly is a compelling alternative.
\end{curiobox}

\vspace{6pt}

\begin{curiobox}{What is the difference between a CDN cache and an in-process cache?}
An \textbf{in-process cache} lives in the same memory space as the application (e.g., a Python dictionary, a Java ConcurrentHashMap, or a Go sync.Map). Reads are in nanoseconds, there is no network hop, and no serialization. The limitation: each application process has its own cache, so processes do not share cached data. If you have 100 application servers, you have 100 separate caches. Cache invalidation requires sending an invalidation signal to all 100 processes. An \textbf{out-of-process cache} (like Redis) is a shared external service. All 100 application servers read from and write to the same cache. Consistency is trivially maintained because there is only one copy. The cost is a network hop (0.1--1ms in the same data center). A \textbf{CDN cache} (Cloudflare, CloudFront, Fastly) sits in front of your entire application, caching HTTP responses at the edge --- geographically close to users. CDN caches serve static assets, API responses, and rendered HTML without any request reaching your origin server. Response latency from a CDN edge is 1--20ms globally; a request to your origin might take 100--500ms. CDN caches are the outermost layer; Redis is the application layer; in-process caches are the innermost layer. Well-designed systems use all three.
\end{curiobox}

\vspace{6pt}

\begin{curiobox}{How does Redis handle the ``O(1) everything'' promise? LRU eviction seems expensive.}
Redis's LRU is approximate, not exact. An exact LRU requires a sorted doubly-linked list of all keys by access time --- maintaining this at millions of operations per second would be prohibitively expensive. Instead, Redis uses a random sampling approach: when it needs to evict a key, it samples a small number of random keys (default: 5), compares their last access times (stored as a 24-bit timestamp in the object header), and evicts the least recently used among the sample. With a sample of 5, this is a good approximation of true LRU. With a sample of 10, it approaches true LRU quality. The \texttt{maxmemory-samples} config controls this. The same sampling approach applies to LFU eviction. This is why Redis's eviction is described as ``approximate LRU'' in the documentation --- it is not exact, but in practice the error is imperceptible.
\end{curiobox}

\vspace{6pt}

\begin{curiobox}{What is cache penetration, and how is it different from cache stampede?}
Cache penetration and cache stampede are both cache failure modes, but they attack different weaknesses. \textbf{Cache stampede} occurs when a hot key expires and many concurrent requests for that key simultaneously miss the cache. The cache \emph{did} have the key; it just expired. \textbf{Cache penetration} occurs when requests are made for keys that \emph{never exist} in the cache or the database. An attacker can deliberately query for millions of random, invalid IDs (user IDs that don't exist, product IDs that were never created). Every query misses the cache (obviously, the key doesn't exist) and hits the database (which finds nothing). The database is hit on every single request for no benefit. The solution is to cache negative results: if the database returns empty for key X, store \texttt{SET missing:X 1 EX 300} and return the empty result from cache for 5 minutes. A second defense is a Bloom filter: a probabilistic data structure that can definitively say ``this key definitely does not exist'' with no false negatives. Place the Bloom filter in front of the cache; if the key fails the Bloom filter, return 404 immediately without querying cache or database.
\end{curiobox}

\vspace{6pt}

\begin{curiobox}{How does Cloudflare use caching at a scale of tens of millions of requests per second?}
Cloudflare operates over 300 Points of Presence (PoPs) globally as of 2024, serving over 50 million HTTP requests per second at peak. Each PoP runs an in-memory cache (based on Nginx's proxy cache, supplemented by custom software). The cache is populated on the first cache miss at a given PoP: the first request for a resource triggers a single origin fetch; all subsequent requests for that resource at that PoP are served from the PoP's local memory without reaching the origin. Cloudflare's custom protocol, called ``Argo,'' additionally routes cache misses through a private backbone network to find the closest PoP that has already cached the resource rather than always going to origin. This ``tiered caching'' reduces origin load by an additional 30--50\% on top of the first-layer PoP caching. The entire Cloudflare cache operates using a variant of stale-while-revalidate: a cached response is served immediately even if its TTL has expired, while a background request refreshes it. The user experiences zero additional latency from the TTL expiry.
\end{curiobox}

\vspace{6pt}

\begin{curiobox}{Can you use a cache for write-heavy workloads, or is caching only for reads?}
Caching is typically described as a read optimization, but it is used for write-heavy workloads in two important ways. First, \textbf{write coalescing}: if 1,000 events update the same counter in one second (e.g., view counts on a viral video), you can buffer those updates in Redis using \texttt{INCR} and periodically flush the accumulated count to the database in a single write. This converts 1,000 database writes into one, at the cost of a few seconds of propagation delay. YouTube and Snapchat use this pattern for view count updates. Second, \textbf{write buffering via write-back}: the cache absorbs writes and asynchronously drains to the database during off-peak periods. The danger, as discussed in the write-back section, is data loss on cache failure. For counters and analytics (where losing a few seconds of data is acceptable), this is a valid tradeoff. For financial data, it is not.
\end{curiobox}

% ── Self-Quiz ─────────────────────────────────────────────────────────────────
\newpage
\section*{Self-Quiz}
\addcontentsline{toc}{section}{Self-Quiz}

\textit{Cover the answers on the next page. Answer each question as if you are in an interview.}

\vspace{8pt}

\begin{quizbox}
\textbf{Q1 [Pattern Selection]:} You are designing a product detail page for an e-commerce platform serving 5 million daily page views. Product data changes infrequently (once a day on average) but is read millions of times. The database query to build a full product page takes 50ms. Walk me through your caching strategy: which write pattern you would use, what TTL you would set, and how you handle cache invalidation when a product's price changes.

\vspace{10pt}

\textbf{Q2 [Failure Mode]:} Your engineering team reports that after a scheduled nightly job that updates 500,000 product records simultaneously, your database CPU spikes to 100\% for 3 minutes and your API latency degrades significantly. No code changed. What is happening and how would you fix it?

\vspace{10pt}

\textbf{Q3 [Data Structure Selection]:} Design the data layer for a real-time gaming leaderboard with 10 million players where: (a) players update their scores frequently during an active game session, (b) you need to serve the global top-100 in under 10ms, (c) you need to tell any player their exact rank in under 10ms. Specify the data structure, the Redis commands, and the time complexity of each operation.

\vspace{10pt}

\textbf{Q4 [Distributed Caching]:} You start with a Redis Cluster of 6 nodes (3 primary, 3 replica). After 6 months, your data grows beyond the cluster's capacity and you add 2 new primary nodes. What happens to your cached data? What is the difference between Redis Cluster's approach to this problem and naive modulo hashing?

\vspace{10pt}

\textbf{Q5 [Consistency vs Performance]:} A payment service caches user account balances in Redis to avoid hitting the PostgreSQL database on every authorization check. A senior engineer raises a concern: ``If a user's balance is updated in PostgreSQL but the Redis cache still has the old value, we might approve a payment the user can't afford.'' How would you design the caching strategy to prevent this? What write pattern would you use and why?
\end{quizbox}

\newpage
\begin{answerbox}
\textbf{A1:} Use cache-aside. On a read miss, query the database (50ms), serialize the result, store it in Redis with \texttt{SET product:\{id\}:page \{json\} EX 3600} (1-hour TTL). On a price change, the application writes to the database first, then calls \texttt{DEL product:\{id\}:page} to invalidate the cache. The next request repopulates it. Set TTL as a safety net even with explicit invalidation --- if the invalidation call fails (Redis briefly unavailable), the entry expires in 1 hour rather than serving stale data indefinitely. Stagger TTLs with ±10\% random jitter to prevent synchronized mass expiry across all products.

\vspace{8pt}

\textbf{A2:} This is a cache stampede triggered by simultaneous TTL expiry. The nightly job updates 500,000 records. If those records were all cached with the same TTL, they were all written to the cache at roughly the same time and they all expire at roughly the same time --- typically around the next nightly job run. When they expire simultaneously, 500,000 cache keys go cold at the same instant, and the full read traffic for those products hits the database concurrently. Fix: (1) Add random TTL jitter --- instead of a fixed 3600-second TTL, use \texttt{3600 + random.randint(-300, 300)} so expirations are spread over a 10-minute window rather than a single instant. (2) Use stale-while-revalidate: serve the stale cached value while a background process refreshes it. (3) After the nightly job writes to the database, proactively rewrite the cache entries (write-through on the job's writes) so the cache is never stale.

\vspace{8pt}

\textbf{A3:} Use a Redis Sorted Set. \texttt{ZADD leaderboard:global score playerid} (O(log N)) on every score update. \texttt{ZRANGE leaderboard:global 0 99 REV WITHSCORES} (O(log N + 100)) to fetch top-100 --- this is sub-millisecond for 10 million entries. \texttt{ZREVRANK leaderboard:global playerid} (O(log N)) to get a player's rank --- also sub-millisecond. The sorted set maintains order automatically on every \texttt{ZADD}; you never need to sort in application code. For frequent score updates during a game session, buffer updates in-game and call \texttt{ZADD} once per game end, not once per point.

\vspace{8pt}

\textbf{A4:} Redis Cluster uses 16,384 fixed hash slots. When you add 2 new primary nodes, Redis Cluster performs \emph{slot migration}: it moves a subset of slots (and the keys in those slots) from existing nodes to the new nodes. Only the keys in the migrated slots are affected. Other keys remain on their existing nodes, serving traffic normally. The migration happens online --- the cluster is not paused. With naive modulo hashing (hash(key) mod N), adding 2 nodes changes N from 3 to 5, which changes the mapping for $(1 - \frac{3}{5}) = 40\%$ of all keys immediately. Virtually the entire cache goes cold. Consistent hashing (and Redis Cluster's slot-based equivalent) causes only $\frac{2}{5}$ of existing keys to move to the new nodes --- the rest stay where they were.

\vspace{8pt}

\textbf{A5:} For a payment authorization cache, use \textbf{write-through} combined with cache invalidation on balance change. Every time the balance changes in PostgreSQL (debit, credit, refund), the application writes the new balance to both PostgreSQL and Redis synchronously before returning success to the caller. Redis always reflects the latest committed balance. Additionally, set a short TTL (30 seconds) as a backstop against any edge case where the cache and database diverge due to a failed write-through. A failed write-through should be treated as a transaction failure: if the Redis write fails, do not commit the database write and return an error. Cache-aside is wrong here because there is a window between a database write and cache invalidation where old data could be served. Write-back is completely wrong: if the cache contains the authoritative balance and fails before flushing, you have lost financial state.
\end{answerbox}

% ── What's Next ───────────────────────────────────────────────────────────────
\vspace{16pt}
\begin{insightbox}
\textbf{Next: L1-08 --- Message Queues \& Streaming (Kafka, SQS)}

Caching solves the read problem: you make reads fast by keeping data close to the consumer. Message queues solve the write problem: you make writes resilient by decoupling the producer from the consumer. In any notification system (Gate 1's target), the path from ``event occurs'' to ``APNs/FCM receives the delivery request'' runs through a message queue. The queue absorbs traffic spikes, enables retry semantics, provides at-least-once delivery guarantees, and decouples your application servers from your push delivery workers.

L1-08 covers Kafka, SQS, RabbitMQ, and the fundamental concepts of message delivery semantics --- at-most-once, at-least-once, and exactly-once. After L1-08, you will have all three Gate 1 prerequisites complete (L1-01, L1-02, L1-08) and will be ready to reattempt session\_001 with a full, architecturally correct notification system design.

\textbf{Caching connection to L1-08:} A common pattern is to combine caching and queuing. An event (order placed, score updated) is published to Kafka. A consumer reads from Kafka, queries the database, and writes to the cache. This ``cache population worker'' pattern separates cache updates from request handling, eliminating cache misses on hot paths and enabling write coalescing for high-throughput events.
\end{insightbox}

\vspace{12pt}
{\color{highlight!50}\rule{\linewidth}{0.6pt}}
\begin{center}
\small\color{primary!60}
\textit{L1-07 complete} \quad\textbullet\quad
\textit{Mark \texttt{curriculum/progress.md} as done} \quad\textbullet\quad
\textit{One lesson remaining for Gate 1: L1-08 Message Queues}
\end{center}

\end{document}
